\documentclass[11pt]{article}
\usepackage{amsmath,amsthm,amssymb}
\usepackage[margin=1in]{geometry}
\usepackage{booktabs}
\usepackage{array}
\usepackage{stmaryrd}
\usepackage[colorlinks=true,linkcolor=blue,citecolor=blue]{hyperref}

\newtheorem{theorem}{Theorem}
\newtheorem{proposition}[theorem]{Proposition}
\newtheorem{corollary}[theorem]{Corollary}
\newtheorem{lemma}[theorem]{Lemma}
\theoremstyle{definition}
\newtheorem{definition}[theorem]{Definition}
\newtheorem{remark}[theorem]{Remark}
\newtheorem{example}[theorem]{Example}
\newtheorem{problem}[theorem]{Problem}

\newcommand{\SPoly}{\mathsf{SPoly}}
\newcommand{\PolyN}{\mathsf{Poly}_{\mathbb{N}}}
\newcommand{\Poly}{\mathsf{Poly}}
\newcommand{\Cont}{\mathsf{Cont}}
\newcommand{\DCont}{\mathsf{DCont}}
\newcommand{\Cat}{\mathsf{Cat}}
\newcommand{\Set}{\mathsf{Set}}
\newcommand{\Fam}{\mathsf{Fam}}
\newcommand{\N}{\mathbb{N}}
\newcommand{\DD}{\mathbb{D}}
\newcommand{\sem}[1]{\llbracket #1 \rrbracket}
\newcommand{\id}{\mathrm{id}}
\newcommand{\Ob}{\operatorname{Ob}}
\newcommand{\Out}{\operatorname{Out}}
\newcommand{\cod}{\operatorname{cod}}
\newcommand{\im}{\operatorname{im}}
\newcommand{\Diff}{\mathrm{Diff}}
\DeclareMathOperator{\dom}{dom}

\title{The container derivative is a tangent structure:\\
\large a fibrational positioning}
\author{MacBeth\\ \small Kodamai}
\date{3 October 2026}

\begin{document}
\maketitle

\begin{abstract}
The derivative of a container is usually sold as a pun: $\partial(y^n)=n\cdot y^{n-1}$ is the
type of one-hole contexts, and the resemblance to calculus is left as a wink. We remove the
wink. We show that the category $\SPoly$ of (finitary, multivariable) polynomial functors is a
Cartesian differential category whose differential combinator is exactly the
Abbott--Altenkirch--Ghani--McBride one-hole-context derivative $\partial$; consequently
$\partial$ is the vertical part of a genuine Cartesian tangent structure, realised on the nose
by the dual-number substitution $T(p)=p(y+\varepsilon v)\bmod\varepsilon^2$. The proof is a
reflection along the faithful counting functor into the classical Cartesian differential
category of $\mathbb{N}$-polynomials, so every differential axiom is inherited rather than
re-checked. We isolate the one place this picture is routinely misread: the non-trivial content
lives at the level of \emph{maps}, not objects --- a map $1\to1$ is differential-linear iff it is
a linear container $S\cdot y$, so $y^2$ is not linear and the tangent functor is \emph{not} the
trivial $T(p)=p\times p$. We then position the result against the fibrational theory of
first-order differential structures, flag the orientation hazard that the container bifibration
lives on the \emph{reverse}, lens side of that theory, and record a first restriction result:
$T$ is a strict $\lhd$-monoidal base-change functor from directed containers over $k$ to directed
containers over the dual numbers, so the tangent bundle of a small category lives over
$k[\varepsilon]/\varepsilon^2$.
\end{abstract}

\section{Introduction}\label{sec:intro}

\paragraph{The pun, and what is under it.}
In the theory of containers one learns early that data types have derivatives. If a container's
extension is the polynomial functor $y^n$ --- ``$n$-tuples'' --- then its one-hole contexts form
the functor $n\cdot y^{n-1}$: choose which of the $n$ slots is the hole, and the remaining $n-1$
slots still carry data. Abbott, Altenkirch, Ghani and McBride made this precise for all
containers \cite{AAGM2003}: the one-hole-context construction $\partial$ acts on a container
$S\lhd P$ by, for each shape, replacing one position by a hole, and it satisfies the sum,
product and chain rules of elementary calculus. The notation $\partial$ is suggestive, and the
suggestion is usually left as exactly that --- a mnemonic that the Leibniz rule
$\partial(F\times G)=\partial F\times G + F\times\partial G$ happens to hold for data types.

This note argues that the mnemonic is a theorem. The resemblance to calculus is not a
coincidence of low-degree arithmetic; it is the shadow of a \emph{Cartesian differential
category} structure \cite{BCS2009} in which $\partial$ is literally the differential combinator.
Once that is established, a second structure comes for free: every Cartesian differential category
is a Cartesian tangent category \cite{CockettCruttwell2014}, so $\partial$ is the vertical
(fibre-direction) part of a genuine tangent bundle functor $T$, and $T$ is computed by the
oldest trick in differential algebra --- substitute a dual number and read off the
$\varepsilon$-coefficient.

\paragraph{The gap.}
Two bodies of work sit on either side of this statement without quite touching. On one side,
$\partial$ itself: defined and shown to obey the calculus rules \cite{AAGM2003}, with a chain
rule under composition of containers proved explicitly \cite{MacBethChainRule}, but never (to our
knowledge) organised into the axioms of a differential category. On the other side, the
fibrational theory of first-order differential structures of
Capucci--Cruttwell--Ghani--Zanasi \cite{CCGZ2024}, which packages Cartesian differential
categories, tangent categories and their reverse variants uniformly as sections of a fibration ---
but whose running examples are manifolds and lenses, not containers, even though it cites the
containers literature. The missing sentence is the one that says: here is a differential category
whose combinator is the container derivative, and here is where it sits in the fibrational
picture.

\paragraph{The result.}
Let $\SPoly$ be the category of finitary multivariable polynomial functors up to isomorphism,
with substitution as composition (Definition~\ref{def:spoly}). Our main theorem is the following.

\begin{theorem}[Main theorem]\label{thm:main-intro}
$\SPoly$ is a Cartesian differential category whose differential combinator is the AAGM one-hole
derivative $\partial$. Consequently $(\SPoly,\partial)$ is a Cartesian tangent category with
tangent functor $T(n)=2n$, $T(f)=\langle f\circ\pi_1,\ \partial[f]\rangle$, realised by the
dual-number substitution $T(p)=p(y+\varepsilon v)\bmod\varepsilon^2$.
\end{theorem}

The reader meets the precise statement again, with all terms defined, as
Theorem~\ref{thm:main} on page~\pageref{thm:main}.

\paragraph{The method.}
We do not verify the seven differential-category axioms by hand on containers. Instead we exhibit
a \emph{counting functor} $N\colon\SPoly\to\PolyN$ into the classical Cartesian differential
category of polynomials over the rig $\N$, sending a container to the polynomial that counts its
elements (Definition~\ref{def:counting}). We show $N$ preserves all the differential-category
structure --- crucially $N(\partial_iF)=\partial N(F)/\partial m_i$ --- and that $N$ is
\emph{faithful} on isomorphism classes. A faithful functor reflects equalities of parallel maps,
so each axiom, true in $\PolyN$ because polynomials over a rig form a Cartesian differential
category \cite{BCS2009}, descends to $\SPoly$. The whole proof is three short lemmas
(\S\ref{sec:proof}). Reflection is the right tool precisely because the axioms are equations
between parallel morphisms, and that is exactly what faithful functors see.

\paragraph{The correction that makes it non-trivial.}
There is a tempting but vacuous reading of ``differential structure on containers'' at the level of
objects. In a Cartesian differential category one can ask which objects $B$ are \emph{linear},
meaning $TB\cong B\times B$; but $\SPoly$ is a \emph{full} Cartesian differential category, so
every object is linear and the question says nothing. The content is one level down, at
\emph{maps}. We prove (Proposition~\ref{prop:linear}) that a map $1\to1$ is differential-linear
--- equivalently additive, equivalently $\partial[f]=f\circ\pi_2$ --- if and only if it is a linear
container $S\cdot y$. So $y$ is linear but $y^2,y^3,\dots$ are not. This is the certificate that
the tangent functor is genuinely differential and \emph{not} the trivial
$T(p)=p\times p$ (under which every map would be linear). We state the discriminator as a table
(Table~\ref{tab:linear}) because it is the single most useful thing to carry away.

\paragraph{Where it sits, and one hazard.}
The fibrational theory \cite{CCGZ2024} presents a first-order tangent structure as a section of
the \emph{additive codomain fibration}, with reindexing by \emph{pullback}. We record in
\S\ref{sec:orientation} that our dual-number $T$ realises exactly this forward picture --- and we
flag, because the programme it belongs to has been bitten by this three times, that the container
bifibration $\Cont(\cod)=\Fam(\cod^{\mathrm{op}})$ lives on the \emph{opposite}, lens side of the
same theory (the reverse-derivative side, \cite[Ex.~32, Def.~57]{CCGZ2024}). The naive hope ``the
container bifibration \emph{is} the forward tangent fibration'' is therefore false, and false for a
precise reason: the difference between pullback reindexing $\rho^*$ and $(\Sigma_\rho)^{\mathrm{op}}$.
The role of \cite{CCGZ2024} here is a packaging remark that re-delivers the differential category we
have already built; nothing in the proof rests on it.

\paragraph{Does it see categories?}
A directed container is the same thing as a small category \cite{AhmanUustalu2016}; in polynomial
terms, a small category is a $\lhd$-comonoid. So the natural follow-up is whether $T$ respects this
comonoid structure. In \S\ref{sec:directed} we record that it does, in a precise and slightly
surprising sense: $T$ is a \emph{strict $\lhd$-monoidal base-change functor}
$\DCont(k)\to\DCont(\DD)$, where $\DD=k[\varepsilon]/\varepsilon^2$; the monoidal comparison is the
container chain rule, holding on the nose. Hence $T$ sends a small category to a small category ---
but over the dual numbers, not over the original base. This is not a defect of the construction; we
show the obstruction to keeping it over $\Set$ is exactly the second-order term that $\varepsilon^2=0$
sends to the zero morphism. The headline is clean: \emph{the tangent bundle of a category lives over
the dual numbers, a tangent vector at an object $a$ is a morphism out of $a$, and the zero vector is
the identity.}

\paragraph{Outline.}
\S\ref{sec:prelim} fixes containers, the category $\SPoly$, the differential-category axioms we
use, and the counting functor. \S\ref{sec:proof} proves the main theorem by reflection and
unpacks the combinatorial content of the three non-trivial axioms. \S\ref{sec:linear} proves the
map-level linearity discriminator. \S\ref{sec:orientation} positions the result in the fibrational
theory and states the orientation hazard. \S\ref{sec:directed} records the restriction to directed
containers and its honest caveat. \S\ref{sec:conclusion} states novelty honestly and lists the open
edges.

\section{Preliminaries}\label{sec:prelim}

We recall only what the proof uses.

\paragraph{Containers and polynomial functors.}
A \emph{container in $n$ variables} is a set $S$ of \emph{shapes} together with, for each shape
$s\in S$ and each $i\in\{1,\dots,n\}$, a \emph{position set} $P^i_s$. Its extension is the
polynomial functor
\[
  \sem{S\lhd P}\colon\Set^n\to\Set,\qquad
  \vec X\ \longmapsto\ \sum_{s\in S}\ \prod_{i=1}^{n} X_i^{\,P^i_s}.
\]
We work with finitary containers (all $P^i_s$ finite), so only the cardinalities $|P^i_s|$ matter up
to isomorphism. Composition of containers is substitution of one extension into another, written
$G\lhd F$; it corresponds to functor composition \cite{AAGM2003}.

\paragraph{The base category.}
\begin{definition}[$\SPoly$]\label{def:spoly}
Let $\SPoly$ be the category with objects the natural numbers $n\in\N$ (read $n$ as ``$n$
variables'', i.e.\ $\Set^n$); morphisms $n\to m$ the isomorphism classes of $m$-tuples of finitary
containers in $n$ variables (equivalently, the finitary polynomial functors $\Set^n\to\Set^m$);
composition substitution; identities the variable tuples $(y_1,\dots,y_n)$.
\end{definition}

Passing to isomorphism classes makes substitution strictly associative and unital, so $\SPoly$ is a
genuine $1$-category. Finite products are $n\times n'=n+n'$ with the variable-projection tuples as
projections and $0$ terminal. Each hom-set $\SPoly(n,m)$ is a commutative monoid under coproduct of
containers $+$ (unit the empty container $0$), and composition is left-additive,
$(f+g)\circ h=f\circ h+g\circ h$ and $0\circ h=0$, because an outer coproduct distributes over any
substitution on the right. With products and $(+,0)$ computed componentwise and compatibly, $\SPoly$
is a \emph{Cartesian left-additive category} in the sense of \cite[Def.~1.2.1]{BCS2009}.

\paragraph{The differential combinator.}
For $f\colon n\to m$ define $\partial[f]\colon 2n\to m$ by the directional one-hole derivative
\[
  \partial[f](\vec x\,;\vec v)\ :=\ \sum_{i=1}^{n} v_i\cdot\partial_i f,
\]
where $\partial_i f$ is the AAGM one-hole derivative in the $i$-th variable: a shape $s$ with
$k=|P^i_s|$ positions in variable $i$ contributes $k$ shapes, each replacing one of those positions
by a \emph{new} tangent variable $v_i$ and leaving the other $k-1$ as $x_i$. Thus $\partial[f]$ is
the $2n$-variable container that is degree $1$ in the tangent variables. The AAGM endofunctor
$\partial$ is the $v$-linear part of $\partial[-]$, i.e.\ $\partial[-]$ on $\SPoly(1,1)$.

\paragraph{The differential-category axioms we need.}
A \emph{Cartesian differential category} \cite{BCS2009} is a Cartesian left-additive category with a
combinator $D$ taking $f\colon A\to B$ to $D[f]\colon A\times A\to B$ (thought of as $f$'s directional
derivative, linear in the second argument), subject to seven axioms CD.1--CD.7. We will invoke, by
name, only the three with combinatorial content:
\begin{itemize}
  \item \textbf{CD.5 (chain rule)} $D[g\circ f]=D[g]\circ\langle f\circ\pi_1,\ D[f]\rangle$;
  \item \textbf{CD.6} $D[D[f]]\circ\langle\langle a,b\rangle,\langle 0,d\rangle\rangle
        =D[f]\circ\langle a,d\rangle$;
  \item \textbf{CD.7 (symmetry of mixed partials)}
        $D[D[f]]\circ\langle\langle a,b\rangle,\langle c,0\rangle\rangle
        =D[D[f]]\circ\langle\langle a,c\rangle,\langle b,0\rangle\rangle$.
\end{itemize}
A map $f$ is \emph{linear} when $D[f]=f\circ\pi_2$ (its derivative is itself, independent of the
base point); the linear maps are the morphisms along which the differential is additive. Every
Cartesian differential category is a Cartesian tangent category, with tangent functor $T(A)=A\times
A$, $T(f)=\langle f\circ\pi_1, D[f]\rangle$, projection $p=\pi_1$ and zero section $\langle
1,0\rangle$ \cite[Prop.~4.7]{CockettCruttwell2014}.

\paragraph{The counting functor.}
\begin{definition}[$\PolyN$ and the counting functor $N$]\label{def:counting}
Let $\PolyN$ be the Cartesian differential category of polynomials over the rig $\N$: objects $n\in
\N$, morphisms $n\to m$ the $m$-tuples of polynomial functions $\N^n\to\N$ with $\N$-coefficients,
composition substitution, differential the Jacobian directional derivative
$D[p](\vec x;\vec v)=\sum_i v_i\,\partial p/\partial x_i$. (Polynomials over any commutative rig form
a Cartesian differential category --- the motivating example of \cite{BCS2009} --- the axioms using
only $+$, $\times$, $d/dx$ and never subtraction.) Define $N\colon\SPoly\to\PolyN$ on objects by
$N(n)=n$ and on a container $F=S\lhd P$ by its \emph{counting polynomial}
\[
  N(F)(m_1,\dots,m_n)\ =\ \sum_{s\in S}\ \prod_{i=1}^{n} m_i^{\,|P^i_s|}\ \in\ \N[m_1,\dots,m_n].
\]
\end{definition}

\paragraph{Dual numbers and directed containers.}
Write $\DD=k[\varepsilon]/\varepsilon^2$ for the dual numbers over a base rig/ring $k$; these appear
only in \S\ref{sec:directed}. A \emph{directed container} is, by the theorem of
\cite{AhmanUustalu2016}, exactly a small category $C$, presented as the polynomial comonad
\[
  c(y)=\sum_{a\in\Ob C} y^{\Out(a)},\qquad \Out(a)=\{\text{morphisms with domain }a\},
\]
with counit picking $\id_a$ and comultiplication encoding codomain and composition; the comonoids in
$(\Poly,\lhd,y)$ are precisely the small categories.

\section{The theorem and its proof}\label{sec:proof}

\begin{theorem}\label{thm:main}
$\SPoly$ is a Cartesian differential category with differential combinator $D=\partial$. Hence
$(\SPoly,\partial)$ is a Cartesian tangent category with $T(n)=2n$, $T(f)=\langle
f\circ\pi_1,\partial[f]\rangle$, realised by the dual-number substitution $T(p)=p(y+\varepsilon
v)\bmod\varepsilon^2$.
\end{theorem}

The proof is three lemmas: $N$ preserves the structure, $N$ is faithful, faithful functors reflect
the axioms.

\begin{lemma}[$N$ preserves the differential-category data]\label{lem:preserve}
$N$ is a product-preserving, left-additive functor commuting with the differential combinators:
$N(\partial_iF)=\partial N(F)/\partial m_i$, and hence $N(\partial[f])=D[N(f)]$.
\end{lemma}

\begin{proof}
Each clause is a counting identity.
\emph{Composition:} for set-valued polynomial functors the cardinality of a composite is the composite
of cardinalities, $|\sem{G\lhd F}(\vec X)|=|\sem{G}(\sem{F}(\vec X))|$, and $|\sem{G}(-)|$ depends only
on the cardinality of its argument; so $N(g\circ f)=N(g)\circ N(f)$.
\emph{Products, projections, identities:} $\sem{-}$ sends $n\times n'=n+n'$ to the product of counting
rings and variable-projections to coordinate projections.
\emph{Left-additivity:} coproduct of containers $\mapsto$ sum of counting polynomials, $0\mapsto0$.
\emph{Products of functors:} $|F\times G|=|F|\cdot|G|$, so $\times\mapsto\times$.
\emph{Derivative:} $\partial_iF$ replaces one of the $|P^i_s|$ positions in variable $i$ by a tangent
copy, with counting
\[
  \sum_{s}|P^i_s|\,m_i^{\,|P^i_s|-1}\prod_{j\neq i}m_j^{\,|P^j_s|}
  =\frac{\partial}{\partial m_i}N(F);
\]
summing $v_i\cdot(-)$ over $i$ gives $N(\partial[f])=D[N(f)]$.
\end{proof}

\begin{lemma}[$N$ is faithful on isomorphism classes]\label{lem:faithful}
For parallel $f,g\colon n\to m$ in $\SPoly$, $N(f)=N(g)$ implies $f=g$.
\end{lemma}

\begin{proof}
It suffices to treat $m=1$. The monomials of $N(F)=\sum_s\prod_i m_i^{|P^i_s|}$ are $\prod_i
m_i^{k_i}$, and the coefficient of $\prod_i m_i^{k_i}$ counts the shapes $s$ with exponent vector
$(|P^1_s|,\dots,|P^n_s|)=(k_1,\dots,k_n)$. So the formal polynomial $N(F)$ determines the multiset of
exponent vectors, which \emph{is} the isomorphism class of $F$ (shapes up to bijection, positions up
to cardinality). Over the infinite rig $\N$ a polynomial function determines its formal polynomial, so
$N$ is injective on isomorphism classes.
\end{proof}

\begin{lemma}[Reflection]\label{lem:reflect}
Each differential-category axiom holds in $\SPoly$.
\end{lemma}

\begin{proof}
Fix an axiom: an equation $L=R$ of parallel morphisms built from a map $f$ together with
projections, tupling, $+$, $0$, $\times$, $\circ$ and $\partial[-]$. By Lemma~\ref{lem:preserve} $N$
preserves every one of these operations, so $N(L)$ and $N(R)$ are the two sides of the \emph{same}
axiom instantiated in $\PolyN$ at $N(f)$. Since $\PolyN$ is a Cartesian differential category
(Definition~\ref{def:counting}), $N(L)=N(R)$. By Lemma~\ref{lem:faithful} $N$ is faithful, and a
faithful functor reflects equalities of parallel morphisms; hence $L=R$ in $\SPoly$.
\end{proof}

This proves the differential-category claim of Theorem~\ref{thm:main}. The tangent structure is the
canonical one attached to any Cartesian differential category \cite[Prop.~4.7]{CockettCruttwell2014};
under $N$ it reads $N(f)(x+\varepsilon v)=N(f)(x)+\varepsilon\sum_i v_i\partial_i N(f)\bmod\varepsilon^2$,
i.e.\ the dual-number substitution, which Lemmas~\ref{lem:preserve}--\ref{lem:reflect} show computes
$T$ on the nose in $\SPoly$.

\begin{example}[The worked monomial]\label{ex:yn}
For $p=y^n$ the dual-number substitution gives
$(y+\varepsilon v)^n=y^n+\varepsilon\,n v\,y^{n-1}\bmod\varepsilon^2$, so the $\varepsilon$-coefficient
is $n\cdot y^{n-1}=\partial(y^n)$: the one-hole contexts of an $n$-tuple, recovered as a tangent vector.
The base part $y^n$ is the projection $p\circ\pi_1$. For $n=1,2,3$ this is, respectively, $v$;
$2vy$; $3vy^2$ --- the $n$ ways to be a one-hole context.
\end{example}

\paragraph{What reflection hides.}
Lemma~\ref{lem:reflect} certifies every axiom but suppresses the bijections. The three with content
are worth stating as container isomorphisms, both because they are the mathematics and because they
are the parts a referee will want to see by hand.

\emph{CD.5, the chain rule.} At container level this is
$\partial(G\lhd F)\cong(\partial G\lhd F)\times\partial F$, proved directly as an explicit bijection in
\cite[Thm.~1]{MacBethChainRule}: a hole in the composite is a choice of outer $G$-index together with
a hole in the chosen inner block; the leftover splits as ``some other inner block'' ($\partial G\lhd
F$) and ``this inner block'' ($\partial F$). Lemma~\ref{lem:reflect} re-derives it as a corollary of
the classical chain rule plus faithfulness --- a consistency check on that bijection.

\emph{CD.6.} Write $f=S\lhd P$ in one variable. $D[D[f]]$ has two shape families: \emph{Case B}, where
the second derivative pokes a base position (two distinct holes, carrying both tangent directions),
and \emph{Case T}, where it pokes the existing tangent position (one hole). Setting the second base
slot to $0$ annihilates every Case-B shape, leaving Case T, whose container is exactly $D[f]$ with the
surviving tangent direction; this is the required identity.

\emph{CD.7, Clairaut for containers.} Setting the second tangent slot to $0$ leaves only Case B:
shapes $(s,h_1,h_2)$ with distinct holes $h_1\neq h_2\in P_s$ carrying tangent directions $b,c$. The
swap involution $(s,h_1,h_2)\mapsto(s,h_2,h_1)$ is a bijection of ordered distinct hole-pairs
exchanging the two directions, witnessing symmetry of the mixed second derivative. In several variables
a hole is a pair (variable $i$, position in $P^i_s$) and the same swap applies.

\paragraph{Corroboration.}
The single-variable axioms CD.1--CD.7 were checked symbolically with an \emph{abstract} $f$ (hence for
all single-variable containers, $N$ being faithful), and $N(\partial_iF)=\partial_iN(F)$, the chain
rule, and the CD.7 swap on several thousand random multivariable finite containers. A minimal
Mathlib-free Lean~4 development certifies the core computable facts --- the $\varepsilon$-coefficient
identity of Example~\ref{ex:yn} for $n\le5$, the monomial Leibniz and chain-rule shadows, and the
map-level linearity discriminator of \S\ref{sec:linear} --- with no axioms and no appeal to
\texttt{native\_decide}. These are finite oracles corroborating the reflection, not the reflection
itself.

\section{Linearity: the corrected discriminator}\label{sec:linear}

There is a reading of ``differential structure on containers'' at the level of \emph{objects} that is
tempting and vacuous. One can ask which objects $B$ are linear in the sense $TB\cong B\times B$; but
$\SPoly$ is a \emph{full} Cartesian differential category, so every object is linear and the question
is empty (in the language of \cite[Thm.~69]{CCGZ2024}, $\Diff(T)$ is everything). The non-vacuous
statement is one level down.

\begin{proposition}[Linear maps]\label{prop:linear}
A morphism $f\colon1\to1$ is a linear map of the Cartesian differential category --- equivalently an
additive container, equivalently $\partial[f](x,v)=f(v)$ --- if and only if $f=S\cdot y$ for a set $S$
(a linear container). In particular $y^2,y^3,\dots$ and nonzero constants are \emph{not} linear.
\end{proposition}

\begin{proof}
$f$ is linear iff $D[f]=f\circ\pi_2$. Apply $N$: $v\cdot f'(x)=f(v)$ for all $x,v$ as
$\N$-polynomials. Writing $N(f)=\sum_k a_k x^k$, the left side is $\sum_k a_k k\,v x^{k-1}$ and the
right $\sum_k a_k v^k$; equality forces $a_k=0$ for $k\neq1$, i.e.\ $N(f)=a_1 x$. By
Lemma~\ref{lem:faithful}, $f\cong a_1\cdot y=S\cdot y$ with $|S|=a_1$. Conversely
$D[S\cdot y]=S\cdot v=(S\cdot y)(v)$.
\end{proof}

This is the honest non-vacuity certificate. Under the trivial candidate $T(p)=p\times p$ every map is
linear; Proposition~\ref{prop:linear} shows the actual tangent functor is not that --- $y$ is linear
and $y^2$ is not. Table~\ref{tab:linear} is the summary to carry away.

\begin{table}[h]
\centering
\begin{tabular}{llll}
\toprule
container $f$ & $N(f)$ & $\partial[f](x,v)$ & linear? \\
\midrule
$y$        & $x$     & $v$          & yes \\
$S\cdot y$ & $|S|\,x$ & $|S|\,v$     & yes \\
$y^2$      & $x^2$   & $2vx$        & no \\
$y^3$      & $x^3$   & $3vx^2$      & no \\
constant $c$ & $|c|$ & $0$          & no (unless $c=0$) \\
affine $c+ S\cdot y$ & $|c|+|S|x$ & $|S|v$ & no (unless $c=0$) \\
\bottomrule
\end{tabular}
\caption{The map-level linearity discriminator. Linear $\iff$ $\partial[f](x,v)=f(v)$ $\iff$
$f=S\cdot y$. Machine-checked, zero axioms, in the Lean development.}
\label{tab:linear}
\end{table}

\section{The orientation split: a fibrational positioning}\label{sec:orientation}

The fibrational theory of first-order differential structures \cite{CCGZ2024} gives a common home to
differential categories, tangent categories and their reverse variants. Its master notion: a
first-order differential structure on an object $B$ (of a suitable $2$-category) is a fibration with
fibred biproducts together with a \emph{section} of it; the tangent-category instance
\cite[Thm.~56]{CCGZ2024} takes the fibration to be the \emph{additive codomain fibration} --- the
codomain fibration $\cod\colon B^\to\to B$, whose Cartesian maps are pullbacks, refined so each fibre
object is a commutative monoid in the slice --- and a section $T$ sends each object to a bundle over
it. Their Theorem~74 then shows a first-order Cartesian tangent structure of this kind induces a
first-order Cartesian differential category, with combinator $D_T=t^{-1}\circ Tf\circ t$.

Our dual-number $T$ realises exactly this forward picture: it is a section of the additive codomain
fibration, and the induced combinator $D_T$ is $\partial$. In this sense Theorem~74 \emph{re-delivers}
the differential category of \S\ref{sec:proof}. We stress that this is a packaging remark: nothing in
the proof of Theorem~\ref{thm:main} depends on \cite{CCGZ2024}, and the citation is here to locate the
result in its natural taxonomy, not to carry it.

\paragraph{The hazard.}
There is a seductive stronger claim --- that the container bifibration itself \emph{is} the forward
tangent fibration of \cite{CCGZ2024} --- and it is false. The container logic
$\Cont(\cod)=\Fam(\cod^{\mathrm{op}})$ is built on $\cod^{\mathrm{op}}$, and
\cite[Ex.~32]{CCGZ2024} identifies $\cod^{\mathrm{op}}$ of $\Set$ with dependent lenses, i.e.\
polynomial functors, i.e.\ containers; the \emph{reverse} differential side of that theory uses the
dual of the additive-bundle fibration \cite[Def.~57]{CCGZ2024}. So $\Fam(\cod^{\mathrm{op}})$ sits on
the reverse/lens side, not the forward Thm-56/74 side. The distinction is the difference between
pullback reindexing $\rho^*$ and the opposite direct-image $(\Sigma_\rho)^{\mathrm{op}}$ --- a sign
error at the level of fibrations. We name it explicitly because it is the same orientation hazard that
has mis-oriented adjacent results in this programme more than once; the safe statement is the forward
one, on $\cod$ with $\rho^*$, which is what Theorem~\ref{thm:main} uses.

\begin{remark}
The reverse/lens reading --- ``is $\partial$ a reverse-derivative tangent structure on
$\Fam(\cod^{\mathrm{op}})$?'' --- is a genuine and, as far as we know, open question. We do not
address it here; see \S\ref{sec:conclusion}.
\end{remark}

\section{Does the tangent structure see categories?}\label{sec:directed}

A directed container is a small category, i.e.\ a $\lhd$-comonoid \cite{AhmanUustalu2016}. If $T$
respects that comonoid structure it transports small categories to small categories. It does, with one
twist that is the whole point.

\begin{theorem}[Restriction to directed containers]\label{thm:directed}
$T$ is a strict $\lhd$-monoidal functor
\[
  T\colon(\Poly/k,\ \lhd,\ y)\ \longrightarrow\ (\Poly/\DD,\ \lhd_\DD,\ y+\varepsilon v),
  \qquad \DD=k[\varepsilon]/\varepsilon^2,
\]
where $\lhd_\DD$ is substitution after base change along $k\to\DD$. The monoidal comparison is the
container chain rule $\partial(G\lhd F)\cong(\partial G\lhd F)\times\partial F$, which holds on the
nose (laxator the identity). Granting the standard base-change $2$-functoriality of polynomial
composition \cite{GambinoKock2013}, $T$ therefore preserves $\lhd$-comonoids, i.e.\ sends a small
category internal to $k$ to a small category internal to $\DD$: a base-change functor
$\DCont(k)\to\DCont(\DD)$.
\end{theorem}

The content is a single identity. Writing $T(c)(y,v)=c(y)+\varepsilon v\,c'(y)$, define $\lhd_\DD$ by
substituting the dual-number output of $T(d)$ into $c$; then $T(c\lhd d)$ has tangent part
$v\,(c\circ d)'(y)$, and the ordinary chain rule $(c\circ d)'=c'(d)\cdot d'$ is exactly the
$\lhd_\DD$ tangent part. So $T$ is strict $\lhd$-monoidal, and the AAGM container chain rule
(\S\ref{sec:proof}) is this identity read off in the (shape $\lhd$ position) presentation. Comonoid
preservation then follows because strict monoidal functors preserve comonoids --- the one step we do
not re-derive on container morphisms, borrowing instead the base-change $2$-functoriality of
\cite{GambinoKock2013}.

\begin{remark}[Why the base is $\DD$ and not $k$ --- the honest caveat]\label{rem:caveat}
$T$ is \emph{not} an endofunctor of $\DCont(k)$. Over $\Set$ the polynomial $T(c)=c+\varepsilon v\,
\partial c$ carries no comonad structure restricting to $C$: composition is not even total. Present
$T(c)$ as a would-be two-sorted directed container, with base objects $B(a)$ (a faithful copy of $C$)
and tangent objects $Tn(a,f)$ indexed by the morphisms $f$ of $C$, the derivative $\partial c$ giving
each $Tn(a,f)$ the out-positions $\{\star\}\cup(\Out(a)\setminus\{f\})$. The counit forces
$\id_{Tn(a,f)}=\star$, and the sort constraints force $\cod$ and composition --- except that the
composite $\mathrm{comp}_C(g,r)$ of a surviving position $g\in\Out(a)\setminus\{f\}$ with a base
position $r$ can equal the \emph{removed} position $f$, in which case the required composite has no
target among $Tn(a,f)$'s positions. For any category with a non-identity morphism $f\neq\id_a$ the
pair $g=\id_a,r=f$ exhibits this, so the $\Set$ reading fails universally outside the
discrete/terminal cases (verified by complete enumeration on the walking arrow, $\mathbb{Z}/2$, and an
idempotent monoid). The resolution is the dual numbers: that missing composite re-enters the
differentiated direction, producing an $\varepsilon^2$ term, which $\varepsilon^2=0$ sends to the
\emph{zero morphism} --- not a position any $\Set$-container can carry, but a legitimate morphism of a
$\DD$-module-enriched category. This is exactly why $\lhd_\DD$ (base change), and not free
two-variable substitution, is the correct composition.
\end{remark}

The geometric payoff survives the caveat and is clean. The projection $p\colon T(C)\to C$ (base change
$\DD\to k$, $\varepsilon\mapsto0$) has fibre over $a$ the out-morphisms $\Out_C(a)$, and the zero
section $0\colon C\to T(C)$ lands on the identity direction. So:

\begin{quote}
\emph{A tangent vector of a category at an object $a$ is a morphism out of $a$; the zero tangent
vector is the identity $\id_a$; and the tangent bundle lives over the dual numbers.}
\end{quote}

\begin{problem}[open]\label{prob:closed-form}
Give the closed-form multiplication of $T(C)$ as a $\DD$-linear category, via
$\delta_{T(c)}=\mu^{-1}\circ T(\delta_c)$ with $T$ defined on container \emph{morphisms} over $\DD$ (so
the $\varepsilon$-zero morphism is representable), and identify it with a named construction --- a
tangent algebroid / first-order jet category, or $C$ glued to its arrow-bimodule by the
$\varepsilon$-action. The object counts already exclude the obvious $\Set$ candidates (for the walking
arrow, the arrow category has $3$ objects, $C\otimes\DD$ has $4$, and $T(C)$ has $5$).
\end{problem}

\section{Honest novelty, and the open edges}\label{sec:conclusion}

\paragraph{What is and is not new.}
The derivative $\partial$ is \cite{AAGM2003}; the chain rule under $\lhd$ is \cite{MacBethChainRule}.
That set-valued polynomial functors assemble into a Cartesian differential category with $\partial$ as
combinator is, to our knowledge, not recorded in exactly this form, but it is close to folklore: it is
the $\Set$-coefficient instance of ``polynomials over a rig form a Cartesian differential category''
\cite{BCS2009}, and it is adjacent to the differential structure of Joyal's species and analytic
functors. We do not claim the differential category as deep. The contribution is the clean statement,
the reflection proof (which makes the seven axioms inherited rather than re-checked), the resolution of
the object-versus-map confusion (Proposition~\ref{prop:linear}), the fibrational positioning with its
named orientation hazard (\S\ref{sec:orientation}), and the restriction to directed containers with the
dual-number caveat made precise (\S\ref{sec:directed}).

\paragraph{Open edges, stated honestly.}
\emph{(i) First order only.} Both $\partial$ and the forward theory of \cite{CCGZ2024} are first order;
$\varepsilon^2=0$ truncates the Fa\`a di Bruno tower. Higher and iterated derivatives are out of scope.
\emph{(ii) The reverse/lens reading is unaddressed.} Whether $\partial$ is a reverse-derivative tangent
structure on $\Fam(\cod^{\mathrm{op}})$ --- the side on which the container logic actually lives --- is
open; \S\ref{sec:orientation} only explains why it is a different question.
\emph{(iii) The closed form of $T(C)$} is Problem~\ref{prob:closed-form}: we have the underlying object,
the projection, the zero section and the geometry, but not the named multiplication.
\emph{(iv) An endofunctor tangent structure ``on $\DCont$''} would come, if at all, from the
Weil-algebra-indexed family $W\mapsto\DCont(W)$ with base-change functors (Leung's presentation of
tangent structure as an action of Weil algebras), of which $\DD$ is the first infinitesimal; this is
sketched, not proved.

\paragraph{Why it matters.}
For the differential- and tangent-category reader, the result supplies a concrete, combinatorial model
whose differential combinator is a construction from programming-language theory, with a one-line
discriminator (Table~\ref{tab:linear}) separating it from the trivial tangent functor. For the
container and \texttt{Poly} reader, it upgrades the derivative from mnemonic to structure and shows
precisely how far that structure reaches into the directed (categorical) world --- as far as the dual
numbers, and no further over $\Set$. The headline is portable: \emph{the tangent bundle of a category
lives over the dual numbers, the one-hole derivative is its vertical part, and the obstruction to
keeping it over $\Set$ is the $\varepsilon^2=0$ zero morphism.}

\begin{thebibliography}{9}

\bibitem{AAGM2003}
M.~Abbott, T.~Altenkirch, N.~Ghani, C.~McBride,
\emph{Derivatives of containers},
Typed Lambda Calculi and Applications (TLCA 2003), LNCS 2701, Springer, 2003, 16--30.

\bibitem{AhmanUustalu2016}
D.~Ahman, T.~Uustalu,
\emph{Directed containers as categories},
Proc.\ MSFP 2016, EPTCS 207, 2016, 89--98.

\bibitem{BCS2009}
R.~F.~Blute, J.~R.~B.~Cockett, R.~A.~G.~Seely,
\emph{Cartesian differential categories},
Theory and Applications of Categories \textbf{22} (2009), no.~23, 622--672.

\bibitem{CCGZ2024}
M.~Capucci, G.~Cruttwell, N.~Ghani, F.~Zanasi,
\emph{A fibrational theory of first order differential structures},
preprint, arXiv:2409.05763 (2024).

\bibitem{CockettCruttwell2014}
J.~R.~B.~Cockett, G.~S.~H.~Cruttwell,
\emph{Differential structure, tangent structure, and SDG},
Applied Categorical Structures \textbf{22} (2014), no.~2, 331--417.

\bibitem{GambinoKock2013}
N.~Gambino, J.~Kock,
\emph{Polynomial functors and polynomial monads},
Math.\ Proc.\ Cambridge Philos.\ Soc.\ \textbf{154} (2013), no.~1, 153--192.

\bibitem{MacBethChainRule}
MacBeth,
\emph{The chain rule for the container derivative} (Kodamai internal note, 2026);
the isomorphism $\partial(G\lhd F)\cong(\partial G\lhd F)\times\partial F$.

\end{thebibliography}

\end{document}
